\begin{figure}[htbp]
\centering
\begin{tikzpicture}[fig, x=1mm, y=1mm]
  \def\cw{27}  % 列宽
  \def\rh{9.5} % 行距
  % 列标题
  \foreach \c/\t in {0/NVIDIA GPU, 1/华为昇腾 NPU, 2/寒武纪 MLU}
    \node[font=\figfont\footnotesize\bfseries] at (\c*30,6) {\t};
  % 行
  \foreach \r/\lab/\a/\b/\cc in {
      0/深度学习框架/PyTorch 等/PyTorch 等/PyTorch 等,
      1/算子库与编译器/{cuBLAS、cuDNN}/{CANN 算子库}/CNNL,
      2/集合通信库/NCCL/HCCL/CNCL,
      3/运行时与驱动/{CUDA Runtime}/{AscendCL 运行时}/{CNRT、CNDrv},
      4/硬件/GPU/NPU/MLU} {
    \node[note, anchor=east] at (-17,-\r*\rh) {\lab};
    \foreach \c/\txt in {0/\a,1/\b,2/\cc} {
      \ifnum\r=0
        \node[blkw, minimum width=\cw mm, minimum height=7mm, inner sep=1pt, font=\figfont\footnotesize] at (\c*30,-\r*\rh) {\txt};
      \else\ifnum\r=4
        \node[blk, minimum width=\cw mm, minimum height=7mm, inner sep=1pt, font=\figfont\footnotesize] at (\c*30,-\r*\rh) {\txt};
      \else
        \node[blka, minimum width=\cw mm, minimum height=7mm, inner sep=1pt, font=\figfont\footnotesize, line width=0.5pt] at (\c*30,-\r*\rh) {\txt};
      \fi\fi
    }
  }
  % 软件栈括号
  \draw[decorate, decoration={brace, amplitude=3pt}, line width=0.6pt]
    (75,-0.5*\rh-0.3*\rh) -- (75,-3.5*\rh+0.1*\rh) node[midway, right=3pt, note, align=left] {计算\\软件栈};
  % 调用方向
  \draw[->, line width=0.7pt, ink] (-42,1) -- (-42,-38) node[midway, left, note, text width=4mm, align=center] {调\\用\\方\\向};
\end{tikzpicture}
\caption{计算软件栈的层次与三种平台的对应组件。一次矩阵乘法调用自上而下依次经过各层；不同平台的组件职责相似，但实现与行为并不完全等价}
\label{fig:stack}
\end{figure}
